\scanpage{18}

\uline{Topos et structures algébriques}

\margin{$\Delta=$ \unsure{Pseudotopos|Pseudo-topos|bons topos} (conditions a) b) c) de Giraud) foncteurs commutant aux $\varinjlim$, et exacts à g.}

1) Topos et structures algébriques définissables par $\varprojlim$ finies.
\[
\Delta_0=\text{cat. avec }\varprojlim\text{ finies}
\]

$T_0$ th. algébrique rel /$\Delta_0$, \deleted{définie}\inserted{représentée} par une petite $C\in\mathrm{Ob}\,\Delta_0$\\
\deleted{\unsure{Th.}}
\[
T_0(E)\simeq\underline{\mathrm{Hom}}_{\Delta_0}(C,E)\qquad\text{si }E\in\mathrm{Ob}\,\Delta_0
\]

Si $E\in\mathrm{Ob}\,\Delta$, on sait qu'on a
\[
T_0(E)=\underline{\mathrm{Hom}}_{\Delta_0}(C,E]\approxeq\underline{\mathrm{Hom}}_\Delta(\widehat C,E)
\left[\simeq\underline{\mathrm{Hom}}_{\mathrm{top}}(E,\widehat C)^\circ\quad\text{si }E\text{ un topos}\right]
\]
Ainsi \deleted{\unsure{$C$ est la $\Delta_0$}}\inserted{le topos $\widehat C$ est} la \deleted{topos} $\Delta_0$-$\Delta$ enveloppe de $C$.

2) \inserted{$\Delta$-}Enveloppe d'une \inserted{petite} catégorie $I$ (Pb. 1)

On prend la $\Delta_0$-enveloppe $\rho_{\Delta_0}(I)\subset\underline{\mathrm{Hom}}(I,\mathrm{Ens})^\circ$ (sous-cat. pleine des foncteurs ``de présentation finie'' -- sous-catégorie pleine engendrée par $I$ via $\varinjlim$ finies -- on peut le faire en deux crans, rajoutant les $\varinjlim$ finies de diagrammes dans $I$, puis prenant facteurs directs.... puis on prend la $\Delta_0$-$\Delta$-enveloppe de $\rho_{\Delta_0}(I)$, savoir $\widehat{\rho_{\Delta_0}(I)}$.

3) \uline{Sous-théorie d'une $\Delta$-théorie} \inserted{(représentée par un topos)} -- et sous-topos (Pb. 2)

Soit $R$ un topos, $\Sigma$ un ensemble de flèches de $R$. On considère
\[
E\longmapsto\underline{\mathrm{Hom}}_{\Delta,\Sigma}(R,E)\qquad E\in\mathrm{Ob}\,\Delta
\]

On sature $\Sigma$ par ST1, ST2, ST3 (SGA 4 I 9.1.7 c)), d'où $\Sigma'$ et on trouve $R\Sigma'^{-1}$, qui est un \uline{topos} \inserted{$R'$}, et on a
\[
\underline{\mathrm{Hom}}_{\Delta,\Sigma}(R,E)\simeq\underline{\mathrm{Hom}}_\Delta(R',E)
\]

% margin note written sideways
\margin{\uline{ST1} Si $w=vu$, \deleted{\unsure{lorsque}} et \deleted{\unsure{l'une}}\inserted{deux} des flèches $u,v,w$ \deleted{sont} dans $\Sigma$, \deleted{\unsure{alors la}} la troisième l'est. \deleted{\unsure{Contient|Contravariance}} isos}

% margin note written sideways
\margin{\uline{ST2} Stabilité par changement de base (donc par $\varprojlim$ finies)}

% margin note written sideways
\margin{\uline{ST3} Nature locale ($\Leftrightarrow$ stabilité par $\varinjlim$)}

On sait que $R'$ est un sous-topos de $R$, et tout sous-topos de $R$ s'obtient ainsi (loc. cit.). Si $R\simeq\widetilde C$, $C$ petit site,

% archive sheet number: 17
